\documentclass[ecta,nameyear,draft]{econsocart}
\input{preamble}
\externaldocument[][nocite]{build/supp}[supp.pdf]
\externaldocument[][nocite]{build/appendix}[appendix.pdf]
\begin{document}
\begin{frontmatter}
\title{Neural Bellman Operators}
\runtitle{Neural Bellman Operators}
\begin{aug}
\author[id=au1,addressref={add1}]{\fnms{Qian}~\snm{QI}\ead[label=e92ECTA]{qiqian@pku.edu.cn}}
\address[id=add1]{Peking University}
\end{aug}
\begin{abstract}
Neural Bellman Operators construct and learn continuation operators for feasible
decisions under a specified recursive economic objective. We derive explicit
regularity of the original action marginal in a constrained investment economy
with continuous innovations and nested average value at risk, including binding
capacities and moving adverse tails. Paid original-risk value enclosures supply
outward derivative labels. Bounded readouts admit exact finite-design
optimization-gap checks; action-path and constrained-stationarity bounds
state the additional information needed to control original policy loss.
A recursive level-and-contrast theorem retains state-dependent critic anchors.
Complete proofs preserve the own-future construction and qualified links to
stochastic control and nonlinear recursive utility. An independent comparison
fixes 56 learned actors across two risk levels, two information budgets,
seven representations and two learning losses, alongside an original-model
online control. A separate
original-model service qualifies all 3,712 final policies. Action-marginal
fitting lowers its node residual in all 28 comparisons but raises the
independently enclosed sample mean original cost in all 28; every joint
95 percent population mean-gain interval contains zero. It is slower in
24 of 28 primary complete cohorts. Conventional controls dominate the recorded
ReLU cost--value scenarios, including all five timing processes and declared
reuse volumes. The global feature-cover certificate remains unexecuted.
The mathematical conditions connect information to feasible policy decisions;
the completed numerical comparison shows that finite-design marginal fitting
alone has not delivered the intended economic or computational advantage.
\end{abstract}
\begin{keyword}
\kwd{Dynamic programming}\kwd{Neural approximation}\kwd{Recursive risk}
\kwd{Policy qualification}\kwd{Computational economics}
\end{keyword}
\end{frontmatter}
\input{new121/introduction121}
\input{new121/core_main121}
\input{new/centeredrisk92-active99}
\input{new110/continuous-main}
\input{new119/foundations_main}
\input{new121/marginal_main}
\input{new121/numerical121}
\input{new121/discussion121}
\bibliographystyle{ecta-fullname}
\begingroup\emergencystretch=1em
\bibliography{references99,references110-additions,references120,references121}
\endgroup
\end{document}
